The computation completes all fifteen sizes through $N=256$ for both
CFTs, with $\ell=4,\ldots,12$. All deficits are positive. The original
checks and extension checks are reported in Tables~\ref{tab:complement-checks}
and \ref{tab:complement-extension-checks}. The narrower window changes
only the selected observations and fitted parameters; all previous
relative-entropy values are preserved.

For Ising $I\mid\sigma$ and $\sigma\mid I$, the free exponents are
$0.2853$ and $0.2853$,
compared with $1/4$, with coefficients
$0.57468$ and $0.57465$,
compared with $B_{ab}=0.46544$.
For $I\mid\varepsilon$ and $\varepsilon\mid I$, the exponents are
$1.9691$ and $2.0140$,
compared with $2$. For $\sigma\mid\varepsilon$ and $\varepsilon\mid\sigma$,
they are $0.2376$ and
$0.3298$, compared with $1/4$.
Their deficit-to-leading ratios at $N=256,\ell=12$ are
$0.9827$ and
$1.1567$, respectively.
This endpoint is outside the fitting window.

For boson $00\mid ss$ and $ss\mid00$, the free exponents are
$0.9757$ and $0.9748$,
compared with $1$. For $00\mid3s,s$ and its reverse, they are
$4.8120$ and $4.8243$,
compared with $5$, with fitted coefficients
$0.12188$ and $0.12456$,
compared with $5\pi/64\simeq0.24544$.
For $ss\mid3s,s$ and its reverse, the exponents are
$1.9532$ and $1.9559$,
compared with $2$.

These are effective powers over $y\leq0.04$. The coefficient and
exponent are correlated, so a small residual alone does not establish
the analytical asymptote. Tables~\ref{tab:ising-complement-fits} and
\ref{tab:boson-complement-fits} separately compare the free fit and the
amplitude fit with imposed analytical exponent. The added $\ell=11,12$
points test their continuation outside this window; they do not change
the fitted coefficients. Increasing $N$ at fixed $\ell$ makes $y$ smaller
without increasing the number of sites in $A$. The displayed discrepancies
do not by themselves distinguish omitted continuum terms from finite
lattice effects. No lattice-correction fitting formula is introduced.
